\chapter{Introdução}
\label{chap:introducao}

Sistemas de treinamento neural separam, em geral, a intenção humana da cadeia
que escolhe dados, precisão, perda, atualização, estado e aceitação. O Princípio
da Interlocução investiga uma arquitetura em que essas escolhas tornam-se
contratos compilados, causalmente vinculados ao trabalho e auditáveis por
receipts. A proposta dialoga com contratos de software
\parencite{meyer1992contract}, verificação de programas
\parencite{hoare1969axiomatic} e código com prova
\parencite{necula1997pcc}, mas sua validade empírica depende dos artefatos e dos
adversários concretos deste estudo.

\begin{claimbox}{Fronteira da contribuição}
O objeto não é substituir kernels numéricos. É tornar explícita e causal a
autoridade sobre por que, quando e sob quais invariantes uma primitive é usada,
e condicionar a publicação ao julgamento produzido a partir de evidência.
\end{claimbox}

\section{Pergunta de pesquisa}

A pergunta principal é se autoridade, auditoria e julgamento compilados podem
preservar observáveis confirmados e, quando há redundância no stack, reduzir
custo líquido depois de incluído o overhead da própria governança.

\section{Contribuições e limites}

O manuscrito final deverá enunciar contribuições sem converter uma instância
single-host em generalização. Em especial, tempo, ciclos e RSS não são energia,
e qualquer dimensão indisponível permanece indisponível.
